
\documentclass{article}
\usepackage{colortbl}
\usepackage{makecell}
\usepackage{multirow}
\usepackage{supertabular}

\begin{document}

\newcounter{utterance}

\twocolumn

{ \footnotesize  \setcounter{utterance}{1}
\setlength{\tabcolsep}{0pt}
\begin{supertabular}{c@{$\;$}|p{.15\linewidth}@{}p{.15\linewidth}p{.15\linewidth}p{.15\linewidth}p{.15\linewidth}p{.15\linewidth}}

    \# & $\;$A & \multicolumn{4}{c}{Game Master} & $\;\:$B\\
    \hline 

    \theutterance \stepcounter{utterance}  

    & & \multicolumn{4}{p{0.6\linewidth}}{\cellcolor[rgb]{0.9,0.9,0.9}{%
	\makecell[{{p{\linewidth}}}]{% 
	  \tt {\tiny [A$\langle$GM]}  
	 Du bist Experte für 2D-Physik-Puzzlespiele. Die Spielumgebung sieht wie folgt aus:\\ \tt \\ \tt Ein blauer Ball kann vom oberen Bildschirmrand an jeder horizontalen Position (x-Koordinate) fallen gelassen werden. Die vertikale Fallposition ist immer bei y = 100 fixiert.\\ \tt Der Ball fällt durch die Schwerkraft und kann von Hindernissen im Level rollen oder abprallen.\\ \tt Ziel ist es, den blauen Ball zu einem grünen Zielball zu führen, der sich irgendwo zwischen den Hindernissen befindet.\\ \tt Die Hindernisse können Wände, Rampen oder andere physikalische Formen sein, die die Bewegung des Balls beeinflussen.\\ \tt \\ \tt So sieht der aktuelle Spielaufbau aus:\\ \tt 1. Der grüne Zielball befindet sich an einer bekannten, festen Position.\\ \tt 2. Die Positionen und Formen der Hindernisse sind festgelegt und bleiben unverändert.\\ \tt \\ \tt Ich stelle dir nun ein Bild als Keyframe für diesen Versuch zur Verfügung. Das Bild enthält folgende Informationen:\\ \tt 1. Die im Spielaufbau vorhandenen Hindernisse und ihre Positionen.\\ \tt 2. Die Position des grünen Zielballs in der Spielumgebung.\\ \tt \\ \tt Schlagen Sie bitte anhand dieses Bildes und der obigen Informationen die beste ganzzahlige x-Koordinate zwischen 100 und 500 vor, auf die der blaue Ball beim nächsten Mal fallen gelassen werden sollte, um die Chance zu maximieren, den grünen Ball zu erreichen und zu berühren.\\ \tt WICHTIG: Antworten Sie nur im Format: FALLEN: <x>, wobei <x> eine Zahl zwischen 100 und 500 ist. 
	  } 
	   } 
	   } 
	 & \\ 
 

    \theutterance \stepcounter{utterance}  

    & \multicolumn{4}{p{0.6\linewidth}}{\cellcolor[rgb]{0.8,1,0.9}{%
	\makecell[{{p{\linewidth}}}]{% 
	  \tt {\tiny [A$\rangle$GM]}  
	 FALLEN: 300 
	  } 
	   } 
	   } 
	 & & \\ 
 

    \theutterance \stepcounter{utterance}  

    & & & \multicolumn{2}{p{0.3\linewidth}}{\cellcolor[rgb]{0.95,0.95,0.95}{%
	\makecell[{{p{\linewidth}}}]{% 
	  \tt {\tiny [GM$|$GM]}  
	 drop_x=300, collisions=15, goal=False, steps=None, dist=226.8 
	  } 
	   } 
	   } 
	 & & \\ 
 

    \theutterance \stepcounter{utterance}  

    & & \multicolumn{4}{p{0.6\linewidth}}{\cellcolor[rgb]{0.9,0.9,0.9}{%
	\makecell[{{p{\linewidth}}}]{% 
	  \tt {\tiny [A$\langle$GM]}  
	 Du bist ein Experte für 2D-Physik-Puzzlespiele. Die Spielumgebung sieht wie folgt aus:\\ \tt \\ \tt Ein blauer Ball kann vom oberen Bildschirmrand an jeder horizontalen Position (x-Koordinate) fallen gelassen werden. Die vertikale Fallposition ist immer bei y = 100 fixiert.\\ \tt Der Ball fällt durch die Schwerkraft und kann von Hindernissen im Level rollen oder abprallen.\\ \tt Ziel ist es, den blauen Ball zu einem grünen Zielball zu bringen, der sich irgendwo zwischen den Hindernissen befindet.\\ \tt Die Hindernisse können Wände, Rampen oder andere physikalische Formen sein, die die Bewegung des Balls beeinflussen.\\ \tt \\ \tt So sieht der aktuelle Spielaufbau aus:\\ \tt 1. Der grüne Zielball befindet sich an einer bekannten, festen Position.\\ \tt 2. Die Positionen und Formen der Hindernisse sind festgelegt und bleiben unverändert.\\ \tt \\ \tt Im letzten Versuch:\\ \tt Der blaue Ball wurde von der Position (x=300, y=100) fallen gelassen. Der Ball bewegte sich durch das Level, prallte und rollte an Hindernissen ab, erreichte oder berührte den grünen Zielball jedoch NICHT.\\ \tt \\ \tt Ich zeige euch nun eine Bildsequenz als Keyframes aus diesem Versuch:\\ \tt 1. Die Startposition des blauen Balls.\\ \tt 2. Bilder, die an jedem Kollisionspunkt aufgenommen wurden, an dem der Ball auf ein Hindernis traf.\\ \tt 3. Die Endposition des Balls nach dem Stoppen.\\ \tt \\ \tt Bitte schlagt basierend auf diesen Bildern und den obigen Informationen die beste ganzzahlige x-Koordinate zwischen 100 und 500 vor, an der der blaue Ball beim nächsten Mal fallen gelassen werden sollte, um die Chance zu maximieren, den grünen Ball zu erreichen und zu berühren.\\ \tt WICHTIG: Antworten Sie **nur** im Format: FALLEN: <x>, wobei <x> eine Zahl zwischen 100 und 500 ist. 
	  } 
	   } 
	   } 
	 & \\ 
 

    \theutterance \stepcounter{utterance}  

    & \multicolumn{4}{p{0.6\linewidth}}{\cellcolor[rgb]{0.8,1,0.9}{%
	\makecell[{{p{\linewidth}}}]{% 
	  \tt {\tiny [A$\rangle$GM]}  
	 FALLEN: 200 
	  } 
	   } 
	   } 
	 & & \\ 
 

    \theutterance \stepcounter{utterance}  

    & & & \multicolumn{2}{p{0.3\linewidth}}{\cellcolor[rgb]{0.95,0.95,0.95}{%
	\makecell[{{p{\linewidth}}}]{% 
	  \tt {\tiny [GM$|$GM]}  
	 drop_x=200, collisions=20, goal=True, steps=762, dist=35.9 
	  } 
	   } 
	   } 
	 & & \\ 
 

\end{supertabular}
}

\end{document}
